\section{Complete reconstruction of the register and incidence isometry}
\label{rev4:ki:section}

The reconstruction in Section~\ref{sec:memory} preserves positive
geometry, exceptional incidence, and the deformation direction within
a single composition. The initial information for this composition is
the register produced by the enriched APP--TRIT--TPK state: the two
arithmetic sheets preserve their operations; the tritic orientation
determines the regime; transport, updating, and nonadic prolongation
preserve the differences between positions and the uniform component.
The linear realizations below act on this already determined register.
Their function is to transport information and make its invariants
explicit.

The fundamental distinction concerns the preserved domain. The
register \(K\) contains twelve ordered coordinates; the tetrad
\(B_K\) preserves four positions selected by incidence; a geometric
direction preserves a line in a representation space. Each operation
has its own image and fiber. The following construction first preserves
the twelve coordinates and then derives the reductions involved in the
geometry. An exact reconstruction is thus available before extracting
a support or performing a dimensional reduction.

\subsection{Gram form of the dodecaphasic design}
\label{rev4:ki:gram-subsection}

We number the positions from one to twelve and denote their set by
\(\Omega=\{1,\ldots,12\}\). The Paley--Witt encoding of the
enriched state determines the set \(\mathcal H\) of its \(132\)
hexads. A hexad is a six-position support of a weight-six vector in
the extended ternary code. The explicit construction of the code and
the proof that every quintuple belongs to a unique hexad are developed
in Theorem~\ref{thm:lf-code} and Proposition~\ref{prop:lf-witt}.
We use precisely this incidence here, with the positions transported
in the same order as the coordinates of the register.

Let \(V=\RR^{12}\) and \(W=\RR^{\mathcal H}\) be equipped with
their Euclidean inner products. The incidence matrix defines the map
\begin{equation}
 \mathcal I:V\longrightarrow W,
 \qquad (\mathcal I k)_H=\sum_{p\in H}k_p,
 \qquad \mathcal I_{H,p}=\mathbf1_{\{p\in H\}}.
 \label{rev4:ki:incidence-map}
\end{equation}
The symbol \(\mathbf1\in V\) denotes the uniform vector with twelve
components. We define
\[
 \mu(k)=\frac1{12}\mathbf1^{\mathsf T}k,\qquad
 P_0=I_{12}-\frac1{12}\mathbf1\mathbf1^{\mathsf T},\qquad
 V_0=\mathbf1^\perp.
\]
The projector \(P_0\) coincides with \(P_{11}\) from
Section~\ref{sec:memory}; the zero subscript indicates the zero-mean
condition in this presentation. The orthogonal decomposition is
\(k=\mu(k)\mathbf1+P_0k\).

\begin{lemma}[Gram identity and incidence normalization]
\label{rev4:ki:gram}
The incidence satisfies
\begin{equation}
 \mathcal I^{\mathsf T}\mathcal I
       =36I_{12}+30\mathbf1\mathbf1^{\mathsf T}.
 \label{rev4:ki:gram-formula}
\end{equation}
Consequently,
\[
 U_W=\frac16\mathcal I\big|_{V_0}:
 V_0\longrightarrow E_{\rm inc}:=\mathcal I(V_0)
\]
is an isometry onto an eleven-dimensional subspace.
\end{lemma}
\begin{proof}
For a fixed point, the number of quintuples containing it is
\(\binom{11}{4}\). Each hexad containing that point includes
\(\binom54\) of them. The uniqueness of the extension of each
quintuple gives \(66\) hexads per point. For a fixed pair of points,
the same argument gives \(\binom{10}{3}/\binom43=30\) hexads per
pair. The \((p,r)\) entry of
\(\mathcal I^{\mathsf T}\mathcal I\) counts exactly the hexads
containing both \(p\) and \(r\). Its diagonal entries are therefore
\(66\), and all other entries are \(30\), proving
\eqref{rev4:ki:gram-formula}.

For \(z\in V_0\), the uniform term vanishes and
\[
 \left\|\frac16\mathcal I z\right\|^2
 =\frac1{36}z^{\mathsf T}
       \mathcal I^{\mathsf T}\mathcal I z
 =\|z\|^2.
\]
Injectivity and the dimension of the image follow from this equality.
The factor six is thus fixed by the Gram form of the design.
\end{proof}

The incidence count determines the scale of the isometry. Sums over
hexads are linear observables of the positions; their inner products
contain the multiplicity with which each position and each pair occur
in the design. Removing the uniform mode separates these two counting
levels and leaves the quadratic factor \(36\). This operation provides
an explicit link between the twelve-position representation and its
representation on hexads.

\subsection{Preservation of the mean and inversion of the twelve coordinates}
\label{rev4:ki:full-subsection}

The uniform component is preserved through an additional scalar
coordinate. We equip \(\RR\oplus E_{\rm inc}\) with the inner
product
\[
 \langle(\mu,y),(\nu,z)\rangle_*=
 12\mu\nu+\langle y,z\rangle.
\]
The weight twelve is the squared norm of \(\mathbf1\), fixed by the
number of positions. We define
\begin{equation}
 F:V\longrightarrow\RR\oplus E_{\rm inc},\qquad
 F(k)=\left(\mu(k),\frac16\mathcal I P_0k\right).
 \label{rev4:ki:F}
\end{equation}

\begin{theorem}[Isometric isomorphism of the complete register]
\label{rev4:ki:full-isometry}
The map \(F\) is a linear isometric isomorphism. Its inverse is
\begin{equation}
 F^{-1}(\mu,y)=\mu\mathbf1+
                  \frac16P_0\mathcal I^{\mathsf T}y,
 \qquad y\in E_{\rm inc}.
 \label{rev4:ki:F-inverse}
\end{equation}
In particular,
\begin{equation}
 12|\mu(k)|^2+
 \left\|\frac16\mathcal I P_0k\right\|^2=\|k\|^2.
 \label{rev4:ki:F-norm}
\end{equation}
\end{theorem}
\begin{proof}
The uniform--centered decomposition is orthogonal, and
Lemma~\ref{rev4:ki:gram} preserves the norm of its second summand.
This proves \eqref{rev4:ki:F-norm} and injectivity. Moreover,
\[
 \frac1{36}P_0\mathcal I^{\mathsf T}\mathcal I P_0=P_0.
\]
Substituting \(F(k)\) into \eqref{rev4:ki:F-inverse} gives
\(\mu(k)\mathbf1+P_0k=k\). Conversely, if
\(y\in E_{\rm inc}\), there is a unique \(z\in V_0\) with
\(y=\mathcal I z/6\). The proposed formula returns
\(\mu\mathbf1+z\), whose image is \((\mu,y)\). This proves
surjectivity onto the stated codomain and both inversion identities.
\end{proof}

The restriction \(y\in E_{\rm inc}\) expresses an effective linear
compatibility. The operator
\begin{equation}
 \Pi_{\rm inc}=\frac1{36}\mathcal I P_0
                         \mathcal I^{\mathsf T}:W\longrightarrow W
 \label{rev4:ki:inc-projector}
\end{equation}
is the orthogonal projector onto \(E_{\rm inc}\). Indeed, it is
symmetric, and its square reduces to itself by
\eqref{rev4:ki:gram-formula}; it acts as the identity on the image
of \(\mathcal I P_0\). Compatibility is therefore checked through
\(\Pi_{\rm inc}y=y\). For an arbitrary datum in \(W\),
reconstruction followed by incidence returns \(\Pi_{\rm inc}y\);
the orthogonal component \((I_W-\Pi_{\rm inc})y\) quantifies the
information external to this representation of the register.

\begin{proposition}[Covariance of reconstruction]
\label{rev4:ki:equivariance}
Let \(g\) be an automorphism of the design, and let \(\rho(g)\)
and \(\sigma(g)\) be its permutation actions on positions and
hexads. Then
\[
 F\rho(g)=(1\oplus\sigma(g))F.
\]
The mean is preserved, and reconstruction commutes with the
simultaneous transport of the positions and their hexads.
\end{proposition}
\begin{proof}
The equivalence \(p\in H\Longleftrightarrow g(p)\in g(H)\)
implies \(\mathcal I\rho(g)=\sigma(g)\mathcal I\).
Permutations fix \(\mathbf1\), preserve the mean, and commute with
\(P_0\). Substituting these identities into \eqref{rev4:ki:F}
proves the claim. Covariance of the inverse follows by composition
with \(F^{-1}\).
\end{proof}

Isometric normalization and arithmetic inversion serve different
functions. In the Hadamard coordinatization of the register,
\(U=\mathcal H_{12}K\) and
\(K=\mathcal H_{12}U/4\), with
\(\mathcal H_{12}^{\mathsf T}=\mathcal H_{12}\) and
\(\mathcal H_{12}^2=4I_{12}\).
The operator \(J_H=\mathcal H_{12}/2\) is orthogonal. Hence
\(FJ_H\) preserves the norm of the normalized signed register
\(U/2\) and represents it in the enlarged incidence space.
Integer inversion through \(\mathcal H_{12}/4\) also preserves
the corresponding congruence sublattice. Both cases transport the
generated register, with the phase order fixed before the
transformation \cite{hmtX}.

\subsection{Exact composition from the two nonadic memories}
\label{rev4:ki:memory-subsection}

We return to \(M\), \(q\), \(m_0\), and \(m_1\) from
\eqref{eq:mem-data}, together with the reconstruction operator
\(L=(M^*M|_{V_0})^{-1}M^*\) from Section~\ref{sec:memory}.
The compatible domain is
\(\mathcal D=\RR\times\operatorname{im}M\), with
\(m=(m_0,m_1)\). Theorem~\ref{thm:memory-inversion} establishes
the reconstruction
\[
 \mathcal R(q,m)=\frac q{12}\mathbf1+Lm,
 \qquad LM=P_0.
\]
Since \(L\) takes values in \(V_0\), composition with incidence
has a direct expression:
\begin{equation}
 \mathcal C=F\mathcal R,\qquad
 \mathcal C(q,m)=\left(\frac q{12},\frac16\mathcal I Lm\right).
 \label{rev4:ki:memory-incidence}
\end{equation}

\begin{corollary}[Equivalence between compatible memory and complete incidence]
\label{rev4:ki:composition}
The map \(\mathcal C\) is an isomorphism between
\(\mathcal D\) and \(\RR\oplus E_{\rm inc}\), with inverse
\begin{equation}
 \mathcal C^{-1}(\mu,y)=
 \left(12\mu,\frac16MP_0\mathcal I^{\mathsf T}y\right).
 \label{rev4:ki:C-inverse}
\end{equation}
Its preserved quadratic form is
\[
 \|\mathcal C(q,m)\|_*^2=
 \frac{|q|^2}{12}+\|Lm\|^2=\|\mathcal R(q,m)\|^2.
\]
\end{corollary}
\begin{proof}
Formula \eqref{rev4:ki:memory-incidence} follows from
\(P_0\mathbf1=0\) and \(P_0L=L\). The inverse is obtained by
applying \(k\mapsto(\mathbf1^{\mathsf T}k,Mk)\) to
\eqref{rev4:ki:F-inverse}, using \(M\mathbf1=0\).
The composition identities follow from \(LM=P_0\) and
\(MLm=m\) for \(m\in\operatorname{im}M\). The norm equality
is \eqref{rev4:ki:F-norm} applied to \(\mathcal R(q,m)\).
\end{proof}

\begin{figure}[htbp]
\centering
\begin{tikzpicture}[>=Latex, every node/.style={font=\small},
  box/.style={draw,rounded corners=2pt,align=center,
              inner sep=6pt,minimum height=12mm}]
 \node[box] (memory) at (0,0) {Charge and memories\\$(q,m_0,m_1)$};
 \node[box] (register) at (4.1,0) {Complete register\\$K$};
 \node[box] (incidence) at (8.2,0) {Mean and incidence\\$(\mu,y)$};
 \node[box] (direction) at (4.1,-1.75)
    {Geometric direction\\$\ell_K=\RR P_3K$};
 \draw[<->] (memory) -- node[above] {$\mathcal R$} (register);
 \draw[<->] (register) -- node[above] {$F$} (incidence);
 \draw[->] (register) -- node[right] {projection and passage to the line}
    (direction);
\end{tikzpicture}
\caption{The charge and the two compatible memories reconstruct the twelve
coordinates. The mean and the incidence component provide another
reversible representation of the same register. The geometric direction
is a subsequent reduction with positive-dimensional fibers.}
\label{rev4:ki:figure}
\end{figure}

This composition also determines the positive partition in
\ref{sec:memory}. If \(q>0\), its lengths are recovered through
\begin{equation}
 d_j=\frac1{12}+\frac{1}{6q}
             \bigl(P_0\mathcal I^{\mathsf T}y\bigr)_j,
 \qquad y=\frac16\mathcal I P_0K.
 \label{rev4:ki:positive-lengths}
\end{equation}
The twelve conditions \(d_j>0\) are the equations of the positive
chart expressed in incidence coordinates. With
\(t_j=\sum_{r\leq j}d_r\), the ordered minors
\(t_j-t_i=\sum_{i<r\leq j}d_r\) are reconstructed from
\((\mu,y)\) or from \((q,m_0,m_1)\). Thus the positivity of
the minors and the isometric preservation of incidence are related
by explicit maps. The former is an inequality in the ordered chart;
the latter is an identity of inner products.

\subsection{Recovery of polarization and scope of the reductions}
\label{rev4:ki:direction-subsection}

The realization of the twelve positions on the cosets of
\(A_5/C_5\) determines the three-dimensional projector
\[
 P_3=\frac3{60}\sum_{a\in A_5}
               \chi_3(a^{-1})\varrho(a).
\]
Here \(\varrho\) is the permutation representation and \(\chi_3\)
the three-dimensional character chosen in the marked chart. Its
representatives, character values, and effective matrix are fixed in
\eqref{eq:leech-projectors}; Lemma~\ref{lem:leech-u} proves
\(P_3=P_3^{\mathsf T}=P_3^2\),
\(\operatorname{rank}P_3=3\) and \(P_3\mathbf1=0\).
In this chart, the register direction is obtained through
\begin{equation}
 u_K=P_3P_0K=P_3Lm
     =\frac16P_3P_0\mathcal I^{\mathsf T}y.
 \label{rev4:ki:direction-reconstruction}
\end{equation}
The equalities follow from the two inverses already proved. The same
direction is therefore recovered from the memories or from the
incidence component; the mean is preserved to reconstruct \(K\),
even though the three-dimensional projector annihilates it.

For the register in \eqref{eq:leech-K}, the exact evaluation
\eqref{eq:leech-u-norm} gives
\[
 \|u_K\|^2=
       \frac{6638585+2275584\sqrt5}{20}>0.
\]
The line \(\ell_K=\RR u_K\) and the transverse complement
\(V_{10}(K)=V_0\cap u_K^\perp\) are thus determined.
Its projector is
\begin{equation}
 P_{10}(K)=P_0-
       \frac{u_Ku_K^{\mathsf T}}{\|u_K\|^2}.
 \label{rev4:ki:rank-ten}
\end{equation}
Indeed, the second summand is a rank-one orthogonal projector
contained in \(V_0\); its product with \(P_0\), in either order,
equals itself. Squaring \eqref{rev4:ki:rank-ten} recovers the same
operator, and its trace is \(11-1=10\). We obtain
\(V_0=\ell_K\oplus^\perp V_{10}(K)\).
This is the reduction that subsequently enters the polarization of
the twelve lattice planes.

Incidence also transports these operators. On \(E_{\rm inc}\), the
operator \(\widehat P_3=U_WP_3U_W^*\) satisfies
\(\widehat P_3U_W=U_WP_3\); the identity follows from
\(U_W^*U_W=I_{V_0}\). Transporting \(P_{10}(K)\) in the same
way preserves its rank and inner product. The correspondence between
the two representations is therefore established on vectors and
operators, as well as on their dimensions.

The need to preserve the different components admits precise checks.
The centered observation of \(k\) coincides with that of
\(k+c\mathbf1\) for every \(c\), whereas their charges differ
by \(12c\). In particular, \(k=0\) and \(k=\mathbf1\) have
the same centered incidence component and different means. The
coordinate \(\mu\) is exactly the information separating these two
registers. Likewise, \(u_{ak+b\mathbf1}=a u_k\) for \(a\ne0\);
the resulting line coincides with \(\ell_k\), even though the
amplitude and mean change. These identities describe the fibers of
the maps on the stated ambient space.

The extraction \(k\mapsto\{j:k_j\geq729\}\) has other fibers.
For the specific register \(K\), both \(K\) and \(K+e_1\)
have support \(\{4,6,9,10\}\), where \(e_1\) is the first
standard basis vector. These are two different ambient vectors with
the same tetrad. This check characterizes the information loss of the
support map; the terminal selection of the register remains fixed by
its prior construction. The flag \(B_K\subset H^-_\alpha\) also
preserves the hexad determined by the signs before the carry, whose
agreement with the incidence selection is proved in
Proposition~\ref{prop:lf-flag}.

The common result of these compositions is an explicit mathematical
continuity. Compatible memory reconstructs the register; its
uniform--centered decomposition determines an isometric incidence
representation; reconstruction of its lengths determines the positive
chart; and three-dimensional projection determines the deformation
direction. The tetrad and flag preserve the incidences used by the
exceptional constructions, whereas \(F\) preserves the entire
dodecaphasic register. The enriched history, in turn, retains the
orientation, path, and prolongation that place this register within
the joint discrete structure of the continuum.
